% ============================================================================
%  abstract.tex — Lkhlassa dyal TP
%  hna kankteb chno derna f had lprojet
% ============================================================================

% --- b lfernasia ---
\chapter*{Résumé}
\addcontentsline{toc}{chapter}{Résumé}

\begin{center}
{\large
Ce rapport présente les solutions de la troisi\`{e}me s\'{e}rie de Travaux Pratiques (TP3) en langage Java. Au cours de ce TP, nous avons mis en \oe{}uvre les concepts fondamentaux de la programmation orient\'{e}e objet tels que l'encapsulation, les m\'{e}thodes d'instance, et l'interaction entre les objets.\\[0.4cm]

Les exercices couvrent une vari\'{e}t\'{e} de classes allant de la mod\'{e}lisation d'entit\'{e}s simples (Point, Rectangle, Compte Bancaire) jusqu'\`{a} des syst\`{e}mes avec une logique m\'{e}tier plus pouss\'{e}e comme la gestion de stocks (Produit) ou des mini-jeux (Jeu Devinette).\\[0.4cm]

\textbf{Mots-clés :} Java, POO, Classes, Objets, M\'{e}thodes, Encapsulation.
\vfill
}
\end{center}

% --- b nenglisia (khellitha lmrra khra ila bghitha) ---
% \chapter*{Abstract}
% \addcontentsline{toc}{chapter}{Abstract}
%
% \begin{center}
% {\large
% Provide a concise summary of your project in English here.\\[0.4cm]
%
% State the problem, methodology, and key findings.\\[0.4cm]
%
% \textbf{Keywords:} Keyword~1, Keyword~2, Keyword~3, Keyword~4.
% \vfill
% }
% \end{center}
